\section{Simulation Framework}

The simulation engine implements the operator system defined in Section~4 using
a schema-driven configuration. All field definitions, operator coefficients,
nonlinearities, and boundary conditions are loaded from a machine-readable
configuration file, ensuring full reproducibility.

Simulations evolve the triadic fields
\((\phi, \vec{V}, R)\) under the operator pipeline, producing field snapshots,
coherence metrics, and resonance-envelope trajectories.

\subsection{Simulation Configuration Alignment}

To ensure transparency and reproducibility, the operator coefficients appearing
in the evolution equations are mapped directly to the configuration parameters
used in the accompanying repository. Table~\ref{tab:config-map} summarizes the
correspondence.

\begin{table}[h]
\centering
\begin{tabular}{lll}
\hline
Symbol & Meaning & Config Key \\
\hline
$D_\phi$ & scalar diffusion & \texttt{diffusion.scalar} \\
$D_V$ & vector diffusion & \texttt{diffusion.vector} \\
$D_R$ & resonance diffusion & \texttt{diffusion.resonance} \\
$\alpha_\phi, \alpha_V, \alpha_R$ & alignment strengths & \texttt{alignment.*} \\
$\beta_\phi, \beta_V, \beta_R$ & coupling strengths & \texttt{coupling.*} \\
$\gamma_\phi, \gamma_V, \gamma_R$ & activation strengths & \texttt{activation.*} \\
$\lambda_\phi, \lambda_V, \lambda_R$ & damping & \texttt{stabilization.*} \\
$R_{\max}$ & resonance saturation & \texttt{resonance.max} \\
$\kappa$ & coherence gain & \texttt{resonance.coherence\_gain} \\
$\phi^\ast$ & target scalar profile & \texttt{targets.scalar} \\
$\vec{V}_{\mathrm{tar}}$ & target vector field & \texttt{targets.vector} \\
\hline
\end{tabular}
\caption{Mapping between operator coefficients and simulation configuration keys.}
\label{tab:config-map}
\end{table}

A complete example configuration corresponding to the operator forms in
Section~4 is reproduced below:

\begin{verbatim}
{
  "fields": {
    "phi": { "initial": "gradient", "bounds": [0, 1] },
    "V":   { "initial": "rotation", "magnitude": 1.0 },
    "R":   { "initial": 0.1, "max": 1.0 }
  },

  "diffusion": {
    "scalar": 0.02,
    "vector": 0.01,
    "resonance": 0.005
  },

  "alignment": {
    "scalar": 0.5,
    "vector": 0.4,
    "resonance": 0.3
  },

  "coupling": {
    "scalar": 0.2,
    "vector": 0.15,
    "resonance": 0.1
  },

  "activation": {
    "scalar": 0.8,
    "vector": 0.6,
    "resonance": 1.2
  },

  "stabilization": {
    "scalar": 0.1,
    "vector": 0.1,
    "resonance": 0.05
  },

  "resonance": {
    "max": 1.0,
    "coherence_gain": 0.7
  },

  "targets": {
    "scalar": "uniform",
    "vector": "boundary-aligned"
  }
}
\end{verbatim}

\subsection{Example Simulation: Resonance Alignment}

A resonance-alignment simulation demonstrates the interaction of the triadic
fields under the operator system. The scalar field begins with a gradient,
the vector field with a rotational pattern, and the resonance envelope with a
small uniform seed. Coherence metrics converge as the fields align.

\begin{figure}[h]
\centering
\fbox{
    \begin{minipage}[c][2in][c]{3in}
        \centering
        \textit{Simulation Output Placeholder}\\[6pt]
        \texttt{(e.g., resonance alignment over time)}\\[6pt]
        \rule{2.5in}{0.5pt}\\
        \rule{2.5in}{0.5pt}\\
        \rule{2.5in}{0.5pt}
    \end{minipage}
}
\caption{Placeholder for resonance-alignment simulation output.}
\label{fig:sim-placeholder}
\end{figure}

\subsection{Simulation Architecture}

\begin{figure}[h]
\centering
\begin{verbatim}
+-----------------------------------------------------------+
|                   Simulation Architecture                 |
+-----------------------------------------------------------+
| 1. Load Schema ------------------------------------------ |
|    - field definitions (phi, V, R)                        |
|    - operator coefficients                                |
|    - boundary conditions                                  |
|    - nonlinearities and thresholds                        |
+-----------------------------------------------------------+
| 2. Initialize Fields -------------------------------------|
|    phi(x,0), V(x,0), R(x,0)                               |
|    from config: gradient, rotation, uniform, noise        |
+-----------------------------------------------------------+
| 3. Operator Pipeline -------------------------------------|
|    For each timestep:                                     |
|      a) Diffusion:        D_phi ∇²phi, D_V ∇²V, D_R ∇²R   |
|      b) Alignment:        α * (targets - fields)          |
|      c) Coupling:         β * cross-field interactions    |
|      d) Activation:       γ * nonlinear terms             |
|      e) Stabilization:    λ * damping                     |
+-----------------------------------------------------------+
| 4. Coherence Metrics -------------------------------------|
|    C(phi, V) = (∇phi · V) / (||∇phi|| ||V|| + ε)          |
|    R dynamics include logistic + coherence gain           |
+-----------------------------------------------------------+
| 5. Output -------------------------------------------------|
|    - field snapshots                                      |
|    - coherence curves                                     |
|    - resonance envelope evolution                         |
+-----------------------------------------------------------+
\end{verbatim}
\caption{Text-based simulation architecture diagram showing the full
operator pipeline and data flow.}
\label{fig:sim-architecture}
\end{figure}

\subsection{Simulation Workflow Summary}

The simulation engine follows a schema-driven workflow in which all field
definitions, operator coefficients, nonlinearities, and boundary conditions
are loaded from a machine-readable configuration. The triadic fields
$(\phi, \vec{V}, R)$ evolve under the operator pipeline shown in
Figure~\ref{fig:sim-architecture}, and coherence metrics are computed at
each timestep. This ensures full reproducibility and alignment between the
mathematical model, the schema taxonomy, and the implementation.
